\subsection{Ngăn xếp công nghệ hạ tầng triển khai}

Các thành phần hạ tầng như cơ sở dữ liệu, công cụ đóng gói, nền tảng triển khai và hệ thống quản lý phiên bản là những yếu tố cần thiết đối với hệ thống thương mại điện tử bán lẻ điện thoại tích hợp AI Agent của nhóm. CĐây là lớp công nghệ bảo đảm dữ liệu được lưu trữ an toàn, môi trường thực thi của các dịch vụ duy trì tính nhất quán, quá trình triển khai diễn ra ổn định và vòng đời phát triển phần mềm được kiểm soát chặt chẽ. 

\subsubsection{Cơ sở dữ liệu}

Cơ sở dữ liệu là thành phần nền tảng, có vai trò lưu trữ, quản lý và truy xuất thông tin phục vụ các nghiệp vụ cốt lõi của ứng dụng. Việc triển khai hợp lý giúp đảm bảo tính toàn vẹn (integrity), nhất quán (consistency) và an toàn dữ liệu (security). Cơ sở dữ liệu cũng đóng vai trò là trục liên kết, hỗ trợ xử lý giao dịch theo thời gian thực và tạo tiền đề cho việc mở rộng quy mô hệ thống trong tương lai.

\textbf{PostgreSQL} được lựa chọn làm hệ quản trị cơ sở dữ liệu chính do đáp ứng tốt yêu cầu về độ tin cậy, tính toàn vẹn và khả năng chịu tải của hệ thống. PostgreSQL là một giải pháp quan hệ mã nguồn mở lâu đời, hỗ trợ các đặc tính ACID, các tính năng SQL nâng cao, khả năng mở rộng schema linh hoạt và tối ưu hóa truy vấn hiệu quả. Nhờ khả năng vận hành ổn định, hiệu năng cao và cộng đồng hỗ trợ mạnh, PostgreSQL phù hợp để lưu trữ các dữ liệu giao dịch quan trọng như thông tin người dùng, đơn hàng và kho vận trong môi trường nhiều dịch vụ (microservices), đồng thời đảm bảo hệ thống duy trì hiệu suất ổn định khi khối lượng dữ liệu tăng trưởng theo thời gian.

\subsubsection{Hạ tầng và triển khai}

Trong một hệ thống hướng Microservice, hạ tầng triển khai đóng vai trò quyết định đối với tính ổn định, khả năng tự động hóa và hiệu năng vận hành. Việc đóng gói (containerization) và điều phối triển khai (orchestration/deployment) giúp đảm bảo môi trường thực thi thống nhất, dễ mở rộng và có khả năng phục hồi khi xảy ra lỗi.

\begin{itemize}
    \item \textbf{Docker: }Docker được sử dụng để đóng gói từng microservice vào một container độc lập, đảm bảo môi trường chạy thống nhất từ quá trình phát triển đến triển khai thực tế. Container hóa giúp loại bỏ các khác biệt về môi trường giữa các máy, hạn chế lỗi phát sinh do phụ thuộc hệ thống, đồng thời hỗ trợ việc mở rộng dịch vụ dễ dàng hơn. Với Docker, mỗi dịch vụ đều có thể được triển khai, cập nhật hoặc phiên bản hóa một cách linh hoạt, tạo nền tảng vững chắc cho quy trình CI/CD và triển khai đa môi trường.
    \item \textbf{Kubernetes (K8s): }Kubernetes được sử dụng để điều phối và quản lý các container ở quy mô lớn, đảm bảo hệ thống hoạt động ổn định và có khả năng mở rộng theo nhu cầu\cite{k8s_orchestration}. K8s cung cấp các cơ chế tự động như cân bằng tải, phân phối tài nguyên, tự động khởi động lại container khi xảy ra lỗi (self-healing) và mở rộng số lượng bản sao dịch vụ khi lưu lượng truy cập tăng cao. Nhờ đó, Kubernetes giúp tối ưu vận hành, giảm thiểu thời gian gián đoạn và hỗ trợ triển khai liên tục, đồng thời tạo điều kiện để hệ thống duy trì khả năng phục vụ ổn định trong môi trường sản xuất 24/7\cite{k8s_autoscaling}.
\end{itemize}

\subsubsection{Quản lí phiên bản}
Trong quá trình phát triển phần mềm, hệ thống quản lý phiên bản là thành phần không thể thiếu nhằm đảm bảo tính đồng nhất của mã nguồn, kiểm soát lịch sử thay đổi và hỗ trợ làm việc nhóm một cách có tổ chức.

\textbf{GitHub} được sử dụng làm kho lưu trữ mã nguồn tập trung và hệ thống kiểm soát phiên bản chính cho toàn bộ dự án. Nền tảng này cho phép theo dõi chi tiết lịch sử thay đổi, hỗ trợ quy trình phối hợp nhóm hiệu quả thông qua các tính năng như Pull Request, Code Review và phân nhánh (branching). Nhờ đó, GitHub đảm bảo tính toàn vẹn của mã nguồn, tăng tính minh bạch trong phát triển và tạo điều kiện triển khai quy trình làm việc có kiểm soát, rõ ràng và chuyên nghiệp.

\subsubsection{Tóm tắt công nghệ}

\begin{table}[H]
\centering
\renewcommand{\arraystretch}{1.4} % dãn dòng
\caption{Tổng hợp các công nghệ sử dụng tại Hạ tầng}
\label{tab:infrastructure_summary}
\begin{tabular}{|p{3cm}|l|p{5.5cm}|} 
\hline
\textbf{Thành phần} & \textbf{Công nghệ} & \textbf{Ghi chú/Vai trò} \\ 
\hline
Cơ sở dữ liệu & PostgreSQL & CSDL quan hệ chính, đảm bảo tính toàn vẹn và bền vững dữ liệu. \\ 
\hline
Containerization & Docker & Đóng gói mỗi Microservice, đảm bảo môi trường chạy nhất quán. \\ 
\hline
Orchestration / Triển khai & Kubernetes (K8s) & Tự động triển khai, mở rộng (scaling) và tự phục hồi hệ thống. \\ 
\hline
Quản lý Phiên bản & GitHub & Quản lý mã nguồn, kiểm soát phiên bản và cộng tác nhóm. \\ 
\hline
\end{tabular}
\end{table}
